\documentclass{article}
\usepackage{mathtools} 
\usepackage{fontspec}
\usepackage[UTF8]{ctex}
\usepackage{amsthm}
\usepackage{mdframed}
\usepackage{xcolor}
\usepackage{amssymb}
\usepackage{amsmath}
\usepackage{hyperref}
\usepackage{mathrsfs}



% 定义新的带灰色背景的说明环境 zremark
\newmdtheoremenv[
  backgroundcolor=gray!10,
  % 边框与背景一致，边框线会消失
  linecolor=gray!10
]{zremark}{注释}

% 通用矩阵命令: \flexmatrix{矩阵名}{元素符号}{行数}{列数}
\newcommand{\flexmatrix}[4]{
  \[
  #1 = \begin{pmatrix}
    #2_{11}     & #2_{12}     & \cdots & #2_{1#4}   \\
    #2_{21}     & #2_{22}     & \cdots & #2_{2#4}   \\
    \vdots      & \vdots      & \ddots & \vdots     \\
    #2_{#31}    & #2_{#32}    & \cdots & #2_{#3#4}
  \end{pmatrix}
  \]
}

% 简化版命令（默认矩阵名为A，元素符号为a）: \quickmatrix{行数}{列数}
\newcommand{\quickmatrix}[2]{\flexmatrix{A}{a}{#1}{#2}}


\begin{document}
\title{6.3}
\author{张志聪}
\maketitle

\section*{1}

验证定义即可
\begin{itemize}
  \item (i)

        对任意$k_1, k_2 \in \mathbb{C}, f_1, f_2, g \in \mathbb{C}[x]_n$，有
        \begin{align*}
          (k_1 f_1 + k_2 f_2, g)
           & = \sum \limits_{k = 1}^n (k_1 f_1 + k_2 f_2)(k) \overline{g(k)}                                           \\
           & = \sum \limits_{k = 1}^n [k_1 f_1(k) + k_2 f_2(k)] \overline{g(k)}                                        \\
           & = \sum \limits_{k = 1}^n k_1 f_1(k)  \overline{g(k)} + \sum \limits_{k = 1}^n k_2 f_2(k)  \overline{g(k)} \\
           & = k_1 \sum \limits_{k = 1}^n f_1(k)  \overline{g(k)} + k_2 \sum \limits_{k = 1}^n f_2(k)  \overline{g(k)} \\
           & = k_1 (f_1, g) + k_2 (f_2, g)
        \end{align*}

  \item (ii)

        \begin{align*}
          \overline{(g, f)}
           & = \overline{\sum \limits_{k = 1}^n g(k) \overline{f(k)}} \\
           & = \sum \limits_{k = 1}^n \overline{g(k) \overline{f(k)}} \\
           & = \sum \limits_{k = 1}^n \overline{g(k)} f(k)            \\
           & = (f, g)
        \end{align*}

  \item (iii)

        对任意$f \in \mathbb{C}[x]_n$，有
        \begin{align*}
          (f, f) & = \sum \limits_{k = 1}^n f(k) \overline{f(k)} \\
                 & = \sum \limits_{k = 1}^n |f(k)|^2             \\
                 & \geq 0
        \end{align*}

        $f = 0$，显然$(f, f) = 0$；

        $(f, f) = 0$，假设$f \neq 0$，$f$是$n - 1$次多项式，故$f(x) = 0$只有$n - 1$个根，
        这与$(f, f) = 0$，$f$有$n$个根矛盾，故$f = 0$.
\end{itemize}

\section*{2}

\begin{itemize}
  \item 必要性;

        $U$是酉矩阵，故
        \begin{align*}
          U \overline{U^T} = E
        \end{align*}
        设$U = (u_{ij})$，把$U$的行向量组写出：
        \begin{align*}
          \mu_i = (u_{i1} , u_{i2}, \cdots, u_{in}) (i = 1, 2, \cdots, n)
        \end{align*}
        把它们看做$\mathbb{C}^n$中向量组，则$U \overline{U^T} = E$等价于
        \begin{align*}
          (\mu_i, \mu_j) = u_{i1} \overline{u_{j1}} + u_{i2} \overline{u_{j2}}
          + \cdots + u_{in} \overline{u_{jn}}= \delta_{ij}
        \end{align*}
        即$\mu_1, \mu_2, \cdots, \mu_n$是$\mathbb{C}^n$中的一组标准正交基。

        如把$U$的列向量组写出
        \begin{align*}
          \mu_i = \begin{bmatrix}
                    u_{1i} \\
                    u_{2i} \\
                    \vdots \\
                    u_{ni}
                  \end{bmatrix}
          (i = 1, 2, \cdots, n)
        \end{align*}
        把它们看做$\mathbb{C}^n$中向量组，则$\overline{U^T} U = E$等价于
        \begin{align*}
          (\mu_i, \mu_j)
           & = \overline{(\mu_j, \mu_i)}                                                                                      \\
           & = \overline{\mu_{1j} \overline{\mu_{1i}} + \mu_{2j} \overline{\mu_{2i}} + \cdots + \mu_{nj} \overline{\mu_{ni}}} \\
           & = \overline{\mu_{1i}} \mu_{1j} + \overline{\mu_{2i}} \mu_{2j} + \cdots + \overline{\mu_{ni}} \mu_{ni}            \\
           & = \delta_{ij}
        \end{align*}
        即$\mu_1, \mu_2, \cdots, \mu_n$是$\mathbb{C}^n$中的一组标准正交基。

  \item 充分性.

        以行向量组为例，于是我们有
        \begin{align*}
          U \overline{U^T} = E
        \end{align*}
        故$U$是酉矩阵。
\end{itemize}

\section*{3}

\begin{itemize}
  \item (1)

        假设$G$不可逆，设$G$的列向量为
        \begin{align*}
          \mu_i = \begin{bmatrix}
                    g_{1i} \\
                    g_{2i} \\
                    \vdots \\
                    g_{ni}
                  \end{bmatrix}
          (i = 1, 2, \cdots, n)
        \end{align*}
        于是我们有
        \begin{align*}
          k_1 \mu_1 + k_2 \mu_2 + \cdots + k_n \mu_n = 0
        \end{align*}
        其中$k_1, k_2, \cdots, k_n$不全为零，故有
        \begin{align*}
          (\sum \limits_{i = 1}^n k_i \epsilon_i, \epsilon_j) = 0, j = 1, 2, \cdots, n
        \end{align*}
        令
        \begin{align*}
          \alpha = \sum \limits_{i = 1}^n k_i \epsilon_i
        \end{align*}
        显然$\alpha \neq 0$，因为对任意$\epsilon_j(j = 1, 2, \cdots, n)$，有
        \begin{align*}
          (\alpha, \epsilon_j) = 0
        \end{align*}
        于是可得
        \begin{align*}
          (\alpha, \alpha) = 0
        \end{align*}
        出现矛盾。
  \item (2)

        对任意$\alpha = (\epsilon_1, \epsilon_2, \cdots, \epsilon_n) X, \beta = (\epsilon_1, \epsilon_2, \cdots, \epsilon_n) Y$，
        利用本题（3），我们有
        \begin{align*}
          (\alpha, \beta) & = X^T G \overline{Y}
        \end{align*}
        因为
        \begin{align*}
          (\alpha, \beta) = \overline{(\beta, \alpha)}
        \end{align*}
        故
        \begin{align*}
          X^T G \overline{Y}
           & = \overline{Y^T G \overline{X}} \\
           & = \overline{Y^T} \overline{G} X
        \end{align*}
        因为$(\alpha, \beta)$是标量，所以转置不影响取值，故
        \begin{align*}
          \overline{Y^T} \overline{G} X
           & = (\overline{Y^T} \overline{G} X)^T \\
           & = X^T \overline{G}^T \overline{Y}
        \end{align*}
        于是可得
        \begin{align*}
          X^T G \overline{Y} & = X^T \overline{G}^T \overline{Y}
        \end{align*}
        所以
        \begin{align*}
          G = \overline{G}^T
        \end{align*}

  \item (3)

        证明略。
\end{itemize}

\section*{4}

提示：
显然$1, x, x^2$是$\mathbb{C}[x]_n$的一组基，然后对这组基进行正交化和单位化即可。

\section*{5}

证明思路和命题2.3类似。

酉变换在标准正交基下的矩阵$U$是酉矩阵，故
\begin{align*}
  U \overline{U^T} = E
\end{align*}
由习题3-(2)可知$\overline{U^T} = U$，所以
\begin{align*}
  |U \overline{U^T}| = |U| \cdot |\overline{U^T}| = |U|^2 = 1
\end{align*}
故$|U| = \pm 1$.

设$\lambda_0$是任意特征值，由习题3-(1)可知$U$可逆，故$\lambda_0 \neq 0$
（因$f(\lambda_0) = |0 \cdot E - U| = |-U| \neq 0$）.齐次线性方程组
$(\lambda_0 - U) X = 0$在$\mathbb{C}^n$内有一非零解向量$X_0$，于是
$U X_0 = \lambda_0 X_0$，两边取复共轭再转置，得$\overline{X_0^T} \overline{U^T} = \overline{\lambda_0} \overline{X_0^T}$，
因为$\overline{U^T} = U$，所以$\overline{X_0^T} U = \overline{\lambda_0} \overline{X_0^T}$，
等式两边右乘$X_0$，利用$UX_0 = \lambda_0 X_0$，得
\begin{align*}
  \overline{X_0^T} U X_0         & = \overline{\lambda_0} \overline{X_0^T} X_0 \\
  \overline{X_0^T} \lambda_0 X_0 & = \overline{\lambda_0} \overline{X_0^T} X_0
\end{align*}
现设
\begin{align*}
  X_0 = \begin{bmatrix}
          x_{1}  \\
          x_{2}  \\
          \vdots \\
          x_{n}
        \end{bmatrix}
  \neq 0
\end{align*}
则
\begin{align*}
  \overline{X_0^T} X_0 & = (\overline{x_1}, \overline{x_2}, \cdots, \overline{x_n}) \begin{bmatrix}
                                                                                      x_{1}  \\
                                                                                      x_{2}  \\
                                                                                      \vdots \\
                                                                                      x_{n}
                                                                                    \end{bmatrix} \\
                       & = |x_1|^2 + |x_2|^2 + \cdots + |x_n|^2                                    \\
                       & \neq 0
\end{align*}
于是立得$\lambda_0 \overline{\lambda_0} = 1$，即$|\lambda_0| = 1$.

\section*{6}

当$\alpha = 0$时，命题显然正确。
现设$\alpha \neq 0$.令
\begin{align*}
  \gamma = \alpha - \frac{(\alpha, \beta)}{(\beta, \beta)} \beta
\end{align*}
由内积的性质可知$(\beta, \beta)$是实数。

根据内积的正定性，有$(\gamma, \gamma) \geq 0$，即
\begin{align*}
  0 \leq (\gamma, \gamma)
   & = (\alpha - \frac{(\alpha, \beta)}{(\beta, \beta)} \beta, \alpha - \frac{(\alpha, \beta)}{(\beta, \beta)} \beta)                                                \\
   & = (\alpha, \alpha) - \frac{\overline{(\alpha, \beta)}}{(\beta, \beta)} (\alpha, \beta)
  - \frac{(\alpha, \beta)}{(\beta, \beta)} (\beta, \alpha) + \frac{(\alpha, \beta)}{(\beta, \beta)} \frac{\overline{(\alpha, \beta)}}{(\beta, \beta)} (\beta, \beta) \\
\end{align*}
因为
\begin{align*}
  \frac{(\alpha, \beta)}{(\beta, \beta)} (\beta, \alpha)
   & = \frac{(\alpha, \beta)}{(\beta, \beta)} \overline{(\alpha, \beta)}
\end{align*}
故
\begin{align*}
  0 & \leq  (\alpha, \alpha) - \frac{\overline{(\alpha, \beta)}}{(\beta, \beta)} (\alpha, \beta)
  - \frac{(\alpha, \beta)}{(\beta, \beta)} (\beta, \alpha) + \frac{(\alpha, \beta)}{(\beta, \beta)} \frac{\overline{(\alpha, \beta)}}{(\beta, \beta)} (\beta, \beta)                                     \\
    & = (\alpha, \alpha) - 2 \frac{\overline{(\alpha, \beta)}}{(\beta, \beta)} (\alpha, \beta) + \frac{(\alpha, \beta)}{(\beta, \beta)} \frac{\overline{(\alpha, \beta)}}{(\beta, \beta)} (\beta, \beta) \\
    & = (\alpha, \alpha) - 2 \frac{|(\alpha, \beta)|^2}{(\beta, \beta)} + \frac{|(\alpha, \beta)|^2}{(\beta, \beta)}                                                                                     \\
    & = (\alpha, \alpha) - \frac{|(\alpha, \beta)|^2}{(\beta, \beta)}
\end{align*}
由$(\alpha, \alpha) - \frac{|(\alpha, \beta)|^2}{(\beta, \beta)} \geq 0$，得
\begin{align*}
  (\alpha, \alpha)                & \geq \frac{|(\alpha, \beta)|^2}{(\beta, \beta)} \\
  (\alpha, \alpha) (\beta, \beta) & \geq |(\alpha, \beta)|^2                        \\
  |\alpha| |\beta|                & \geq |(\alpha, \beta)|
\end{align*}

等式成立当且仅当$(\gamma, \gamma) = 0$，即$\gamma = 0$，由$\gamma = 0$可得：
\begin{align*}
  \alpha = \frac{(\alpha, \beta)}{(\beta, \beta)} \beta
\end{align*}
此时$\alpha, \beta$是线性相关的。

\section*{7}

我们有
\begin{align*}
  |\alpha + \beta|^2
   & = (\alpha + \beta, \alpha + \beta)                                      \\
   & = (\alpha, \alpha) + (\alpha, \beta) + (\beta, \alpha) + (\beta, \beta) \\
\end{align*}
因为
\begin{align*}
  (\alpha, \beta) = \overline{(\beta, \alpha)}
\end{align*}
故
\begin{align*}
  (\alpha, \beta) + (\beta, \alpha) & = 2 \Re [(\alpha, \beta)]
\end{align*}
$\Re [(\alpha, \beta)]$表示复数$(\alpha, \beta)$的实部。

因为复数的实部小于等于复数的模，所以
\begin{align*}
  \Re [(\alpha, \beta)] \leq |(\alpha, \beta)|
\end{align*}
综上，并利用柯西-布尼雅可夫斯基不等式，我们有
\begin{align*}
  |\alpha + \beta|^2
   & = (\alpha, \alpha) + (\alpha, \beta) + (\beta, \alpha) + (\beta, \beta) \\
   & = (\alpha, \alpha)  + 2 \Re [(\alpha, \beta)] + (\beta, \beta)          \\
   & \leq (\alpha, \alpha)  + 2 |(\alpha, \beta)| + (\beta, \beta)           \\
   & \leq (\alpha, \alpha)  + 2 |\alpha| |\beta| + (\beta, \beta)            \\
   & = (|\alpha| + |\beta|)^2
\end{align*}
故
\begin{align*}
  |\alpha + \beta| \leq |\alpha| + |\beta|
\end{align*}

\section*{8}

\begin{itemize}
  \item (1)

        由内积的正定性，我们有
        \begin{align*}
          d(\alpha, \beta)
           & = |\alpha - \beta|                        \\
           & = \sqrt{(\alpha - \beta, \alpha - \beta)} \\
           & \geq 0
        \end{align*}
        且$d(\alpha, \beta) = 0$当且仅当$\alpha - \beta = 0$，即$\alpha = \beta$.

  \item (2)

        我们有
        \begin{align*}
          d(\alpha, \beta)^2
           & = |\alpha - \beta|^2                      \\
           & = (\alpha - \beta, \alpha - \beta)        \\
           & = (-1)^2 (\alpha - \beta, \alpha - \beta) \\
           & = (-(\alpha - \beta), - (\alpha - \beta)) \\
           & = (\beta - \alpha, \beta - \alpha)        \\
           & = d(\beta, \alpha)^2
        \end{align*}
        由距离的非负性可知
        \begin{align*}
          d(\alpha, \beta) = d(\beta, \alpha)
        \end{align*}

  \item (3)

        由习题7，我们有
        \begin{align*}
          d(\alpha, \gamma)
           & = d(\alpha, \gamma - \beta + \beta)      \\
           & = |\alpha - (\gamma - \beta + \beta)|    \\
           & = |\alpha - \beta + \beta - \gamma|      \\
           & \leq |\alpha - \beta| + |\beta - \gamma| \\
           & = d(\alpha, \beta) + d(\beta, \gamma)
        \end{align*}
\end{itemize}


\end{document}
